%% GENERATED by python/paper/build.py from results/ -- do not edit by hand.
%% Rebuild:  .venv\Scripts\python.exe python\paper\build.py --only US_ESC_FrenchAll
%% Source:   results/esc/escExperiments.csv
\begin{table}[!htb]
\centering
\begin{threeparttable}
\caption{Endogenous design and all French characteristics in US}
\label{table:US_ESC:frenchAll}
\renewcommand{\arraystretch}{1.25}
\begin{tabularx}{.9\textwidth}{p{2.6cm}YYYYY}
\toprule
\textbf{Scenario} & \textbf{CRRA} ($\rho$) & $\bm{\theta}$ \textbf{(2020)} & \textbf{Tax rate} & \textbf{Savings rate} & \textbf{Avg. workweek} \\
\midrule 
 & 0.5 & 0.74 & 14.42\% & 15.98\% & 39.40 \\
Baseline & 1.0 & 0.74 & 14.43\% & 15.38\% & 39.40 \\
 & 2.0 & 0.74 & 14.43\% & 14.85\% & 39.40 \\[.5em]\hline\\[-.75em]
 & 0.5 & 0.74 & 13.05\% & $+0.70$ p.p. & 35.03 \\
Exogenous $\theta$ & 1.0 & 0.74 & 13.82\% & $+0.22$ p.p. & 35.25 \\
 & 2.0 & 0.74 & 14.86\% & $-0.09$ p.p. & 35.44 \\[.5em]\hline\\[-.75em]
 & 0.5 & 1.00 & 13.13\% & $-0.14$ p.p. & 35.17 \\
Endogenous $\theta$ & 1.0 & 0.97 & 13.83\% & $-0.07$ p.p. & 35.46 \\
 & 2.0 & 0.49 & 14.99\% & $-0.01$ p.p. & 35.18 \\[.5em]\hline\\[-.75em]
 & 0.5 & 1.00 & 21.29\% & $-1.32$ p.p. & 35.44 \\
France & 1.0 & 1.00 & 21.29\% & $-1.95$ p.p. & 35.44 \\
 & 2.0 & 1.00 & 21.29\% & $-2.38$ p.p. & 35.44 \\[.5em]\hline\\[-.75em]
\bottomrule
\end{tabularx}
\begin{tablenotes}
\footnotesize
\item \textit{Note:} Deadweight-cost specification: the proportional cost $f(\theta)$ with $\phi = 0.5$ and $p$ calibrated per $\rho$ (Table~\ref{table:US_ESC:calibration}). Every counterfactual is a separate equilibrium path: the changed parameters hold throughout, the economy starts from its own steady state, and the political choice binds from the first period of the horizon, so the design in force in 2020 is itself an outcome rather than an inherited datum. All rows are read at 2020. $\theta$ (2020) is the design in force there; in the exogenous rows it is the US design, held fixed so that the counterfactual is about the changed characteristic alone. Each $\rho$ is separately calibrated and its workweek normalised against its own baseline. The savings rate is savings relative to GDP; the baseline rows report its level and every other row the change against the baseline at the same $\rho$, in percentage points. The scenario replaces the US $\eta_i$, the level of $X_i$ and the voting weights $\mu_i$ with France\textquotesingle s simultaneously. Common-$X$ calibration: see the note to Table~\ref{table:US:Calib}. The France row is not a counterfactual on the US model: France carries its own characteristics \emph{and} its own calibrated $\omega$, so the distance between it and the endogenous row is what the observable characteristics do not explain. Its workweek is France's own calibration target, not a prediction, and its savings rate is likewise the distance from the US baseline.
\end{tablenotes}
\end{threeparttable}
\end{table}
